% governance.tex -- Section 8: Meta-Governance Engine

\section{Meta-Governance Engine}
\label{sec:governance}

The meta-governance engine is Kant's most distinctive architectural feature:
it allows players to \emph{change the rules of the game during play}. By
proposing, voting on, and applying rule modifications, agents participate in
a form of constitutional self-governance that directly parallels Constitutional
AI~\citep{bai2022constitutional} but operates through democratic mechanism
design rather than static principles.

\subsection{Constitutional Mechanics}
\label{sec:constitutional}

The governance system operates as a layer atop any N-player game. Between
action rounds, a governance phase allows players to propose modifications to
the game's rules, vote on proposals, and apply approved changes. This creates
a two-level game: the \emph{object game} (the underlying strategic interaction)
and the \emph{meta-game} (the governance process itself). Mutable state is
maintained in a \texttt{RuntimeRules} object that tracks enforcement mode,
penalty parameters, active mechanisms, and governance history.

\subsection{Proposal System}
\label{sec:proposals}

Three types of proposals are supported:

\begin{description}[nosep]
    \item[Parameter Proposals.] Modify numeric game parameters (e.g., payoff
        values, contribution multipliers, round counts) within pre-defined
        bounds.
    \item[Mechanic Proposals.] Activate or deactivate one of six composable
        mechanisms (Section~\ref{sec:mechanisms}).
    \item[Custom Proposals.] Free-form modifications specified as payoff
        delta functions, subject to delta-clamping safety constraints
        (Section~\ref{sec:delta-clamping}).
\end{description}

At most $3$ proposals may be submitted per governance round.

\subsection{Democratic Voting}
\label{sec:voting}

Proposals are decided by strict majority vote. The threshold is:
\[
\tau = \left\lfloor \frac{n_{\text{active}} \cdot \alpha}{\beta} \right\rfloor + 1
\]
with default $\alpha/\beta = 1/2$, so $\tau = \lfloor n/2 \rfloor + 1$
(e.g., $3$ of $4$ players must approve). Each active player casts one
approve/reject vote per proposal. Approved proposals are applied immediately
via \texttt{\_apply\_proposal}, modifying the \texttt{RuntimeRules} state.

\subsection{Composable Mechanisms}
\label{sec:mechanisms}

Six governance mechanisms can be proposed, voted on, and composed.
Mechanisms are applied in fixed order (taxation $\to$ redistribution
$\to$ insurance $\to$ quota $\to$ subsidy $\to$ veto), so cumulative
effects are deterministic.

\paragraph{Taxation.} Rate $r = 1/10$ (default). Pool
$P = \sum_i u_i \cdot r$; share $= P / n$. Each player receives
$\hat{u}_i = u_i(1 - r) + P/n$.

\paragraph{Redistribution.} Equal mode sets $\hat{u}_i = \bar{u}$ for
all $i$. Proportional mode applies damping $d = 1/2$:
$\hat{u}_i = u_i + d(\bar{u} - u_i)$.

\paragraph{Insurance.} Contribution rate $c = 1/10$; threshold
$\theta = \bar{u}/2$. Players below $\theta$ split the pool:
payout $= P / |\{i : u_i < \theta\}|$.

\paragraph{Quota.} Cap $q = 8$. Excess $E = \sum_i \max(0, u_i - q)$
is redistributed equally to below-cap players.

\paragraph{Subsidy.} Floor $f = 2$; fund rate $= 1/5$. Above-floor
players contribute $(u_i - f) \cdot 1/5$; below-floor players receive
proportional payouts up to the floor.

\paragraph{Veto.} A designated veto player (default: player $0$) can
trigger equalization: if $u_{\text{veto}} < \bar{u}$, all payoffs are
set to $\bar{u}$.

\subsection{Custom Modifiers with Delta-Clamping Safety}
\label{sec:delta-clamping}

Custom proposals specify arbitrary payoff modifications via delta functions.
To prevent destabilizing changes, a delta-clamping mechanism bounds the
magnitude of any single modification:
\[
\delta_i^{\text{clamped}} = \text{clamp}\!\bigl(\delta_i,\;
  -\max(|u_i| \cdot \rho,\, \rho),\;
  +\max(|u_i| \cdot \rho,\, \rho)\bigr)
\]
where $\rho = 1/2$ (default). This ensures no custom modifier can change
any player's payoff by more than $50\%$ of its absolute value (or $0.5$
in absolute terms, whichever is larger), preventing adversarial proposals
from destabilizing the game.

\subsection{Governance-Coalition Interaction Loop}
\label{sec:gov-coalition}

When coalition formation and governance are active simultaneously, a rich
interaction loop emerges: coalitions may coordinate voting blocs, propose
self-serving rules, or use governance to enforce coalition agreements. This
creates emergent institutional dynamics analogous to real-world political
systems.

The payoff modification chain is: base game payoff $\to$ coalition penalties
$\to$ side payments $\to$ governance mechanisms (in order) $\to$ custom
modifiers (with delta-clamping).

\subsection{Connection to Constitutional AI}
\label{sec:constitutional-ai}

The governance engine operationalizes key ideas from Constitutional
AI~\citep{bai2022constitutional} in a multi-agent game-theoretic setting.
Rather than encoding principles as static rules, Kant allows principles to
emerge through democratic deliberation. This tests whether agents can
\emph{construct} fair governance systems, not merely follow pre-specified ones.

Where Constitutional AI uses static principles curated by designers,
Kant's governance engine tests whether agents can \emph{discover} fair
rules through democratic deliberation---a stronger test of alignment
that does not presuppose the correct principles.
